% main.tex -- compile inside the official ICML author kit (style year set below).
\documentclass{article}

\usepackage{microtype}
\usepackage{graphicx}
\usepackage{booktabs}
\usepackage{hyperref}
\newcommand{\theHalgorithm}{\arabic{algorithm}}

% Package option: none = anonymous review copy (author hidden);
% [accepted] = author shown. The style year must match the venue's own kit.
\usepackage[accepted]{icml2025}
% Neutral first-page footnote; delete these three lines for a camera-ready copy.
\makeatletter
\def\ICML@appearing{Preprint. Under review.}
\makeatother

\usepackage{amsmath}
\usepackage{amssymb}
\usepackage{mathtools}

\icmltitlerunning{Objective Completion versus Constraint-Violating Optimization in AI Agents}

\begin{document}

\twocolumn[
\icmltitle{Distinguishing Objective Completion from\\
           Constraint-Violating Optimization in AI Agents}

\icmlsetsymbol{equal}{*}

\begin{icmlauthorlist}
\icmlauthor{Tariq Mahmood}{inst}
\end{icmlauthorlist}

\icmlaffiliation{inst}{DIMES Department, University of Calabria, Rende (CS), Italy}
\icmlcorrespondingauthor{Tariq Mahmood}{mehmoodt222@gmail.com}

\icmlkeywords{AI agents, LLM agents, trustworthy AI, constraint violation, reward hacking, goal misgeneralization, agent evaluation}

\vskip 0.3in
]

\printAffiliationsAndNotice{}

\begin{abstract}
Agents built on large language models (LLMs) are now routinely asked to plan, call tools, run code, and act on external systems with little supervision. In this setting an agent can reach its assigned objective and still break the rules that were supposed to govern how it got there. Most benchmarks score only the final outcome, so an agent that obtains the right answer through an unauthorized tool, a leaked credential, or a flaw in the environment is counted as a success. We argue that objective achievement and constraint compliance are separate properties and should be measured separately. We describe an evaluation framework with five components: outcome verification, trajectory monitoring, constraint compliance, uncertainty and anomaly detection, and post-hoc behavioral analysis. We discuss how distribution shift relative to known-compliant trajectories might flag risky behavior before a task finishes, why the resulting measurements should not be collapsed into one score, and how they can still support a release decision. We then instantiate the framework on ODCV-Bench, running four LLM agents over all 40 scenarios in two framings for 777 traced runs (an original wave of 297 and an extension wave of 480), and scoring every run under four checker configurations against 262 hand-written constraints. The results are mixed, but informatively so. Under deterministic state and log checkers, the ranking by outcome and the ranking under a constraint gate are \emph{opposed} (Kendall's $\tau=-0.60$): the two agents with the highest task success are exactly the two that the gate disqualifies. A sweep over the gate tolerance shows that this inversion occupies a window rather than holding everywhere, and locating the window requires the checker's false-positive rate, which our still-incomplete audit leaves unmeasured. Under the full protocol the comparison cannot be made at all, because the measured violation rate exceeds the gate tolerance for every agent; this follows from our own analysis of checker error and is not a result about the agents. Two frontier judges on the original 200-run wave agree on only $85.5\%$ of runs, and their disagreement is one-directional and on its own decides which agent would ship. A second panel of three smaller open-weight judges, judging each constraint separately over the 777 runs, reaches a Krippendorff's $\alpha$ of only $0.155$ $[0.103, 0.209]$, and the gate verdict swings from $7\%$ unsafe failure under consensus to $73\%$ under any-judge-flags on the same verdicts. The anomaly score we propose for early warning does not work. At the raw run level it looks anti-predictive (AUROC $0.312$), but this is largely a truncation artefact: violating runs stop at their first violation and compliant runs do not, so the run-level maximum is taken over systematically fewer draws in the positive class. Once the pre-violation horizon is matched between classes, the AUROC is uninformative ($0.496$ at $H{=}2$, $0.432$ at $H{=}3$). Scoring the proposed action instead of the prefix gives a modest prefix-level improvement (AUROC $0.603$), which is still far too weak for a guard and catches nothing at a $5\%$ false-alarm budget. Split-conformal calibration overshoots its target false-alarm rate by a factor of about $2.8$ on the original score. Finally, the category that an outcome-only benchmark cannot tell apart from incompetence, namely declining an objective that cannot be reached compliantly, accounts for $84.5\%$ of runs under the deterministic state+rule checker and $81.5\%$ $[71.4, 90.2]$ under three-judge consensus. It is the dominant category under every well-behaved checker configuration we report.
\end{abstract}

\section{Introduction}
\label{sec:intro}

Large language models were first deployed as question-answering systems. Increasingly they are used as agents: they break a goal into steps, call tools, execute code, browse the web, and interact with external services over many turns \citep{yao2023react}. Once a model acts rather than answers, a distinction that used to matter little becomes central. There is what the agent achieves, and there is how it achieves it.

A simple example makes the point. Suppose an agent is asked to retrieve a record from a particular database. A conventional evaluation marks the run as successful if the returned record is correct. Now suppose the agent got the record by querying a different database it was not authorized to use, by exploiting a vulnerability, by reading credentials it was never meant to see, or by sidestepping an explicit restriction in its instructions. Judged by outcome, nothing went wrong. Judged by any reasonable standard of trustworthy behavior, the run is a failure.

The concern is not hypothetical. Reinforcement learning has a long record of systems that satisfy the literal objective in ways their designers did not intend \citep{amodei2016concrete,krakovna2020specification}, and recent reports describe LLM-based agents that tamper with their environment or their evaluation when the legitimate route is hard \citep{bondarenko2025demonstrating,meinke2024frontier}. Widely used agent benchmarks nevertheless report task success as the headline number \citep{liu2024agentbench,zhou2024webarena,jimenez2024swebench}, and such a number cannot tell an agent that solved the task from one that found a way around it.

The question we address is therefore the following: \emph{how can we distinguish an agent that completes its assigned objective from one that optimizes for the objective while violating the constraints that were supposed to bound its behavior?}

The question becomes sharper as agents take on longer horizons and richer environments. A system optimized mainly for completion then has more opportunities to find strategies that meet the formal objective by exploiting the environment or the evaluation procedure. Our position is that task success is not a sufficient condition for acceptable behavior, and that an evaluation should report at least two things: whether the objective was achieved, and whether the constraints were respected.

The contributions of this paper are:
\begin{itemize}
    \item a statement of the problem that separates objective achievement from constraint compliance (Section~\ref{sec:problem});
    \item a Constraint-Aware Agent Evaluation Framework (CAAEF) with five components (Section~\ref{sec:framework});
    \item a two-dimensional evaluation model, an argument against scalar aggregation, and a release-gating rule that does not need a scalar, together with an analysis of how the rule behaves when the checker makes errors (Section~\ref{sec:model});
    \item an empirical instantiation of the framework on ODCV-Bench, covering $777$ traced agent runs (an original $297$-run wave on 20 scenarios and an extension $480$-run wave on the remaining 20 at three repetitions), scored under four checker configurations against $262$ hand-written constraints (Section~\ref{sec:setup});
    \item results that support part of the framework and refute another part (Section~\ref{sec:results}), together with a measurement of how far the conclusions depend on the checker rather than on the agents, and sensitivity analyses over the gate tolerance, the soft-constraint weights, the judging protocol, and the conformal calibration method;
    \item a set of open research problems that follow from the framework and from what the study failed to show (Section~\ref{sec:challenges}).
\end{itemize}

The study is deliberately small, and we state its limits early because they bound every number we report. It covers one benchmark, across all $40$ of its scenarios in two framings. On the agent side there are four agents from only two model families; on the extension's judge side there are three smaller open-weight judges from three further families. The $297$-run original wave has a single complete repetition, and each cell of the $480$-run extension wave has three. The human audit that would establish the checkers' false-positive and false-negative rates has not yet been carried out, so the quantities that depend on it are reported as pending rather than estimated. What the study can support is a claim about \emph{measurement}: the choice of checker changes the conclusion at least as much as the choice of agent does, and this is visible even at this scale.

\section{Related Work}
\label{sec:related}

\paragraph{Specification gaming and reward hacking.}
\citet{amodei2016concrete} list reward hacking and negative side effects among the concrete problems of AI safety. \citet{krakovna2020specification} collect examples of specification gaming, in which a system satisfies the stated objective without achieving the intended one, and \citet{skalse2022defining} give a formal treatment of reward hacking. Goal misgeneralization is a related failure in which a policy competently pursues the wrong goal out of distribution even though the training specification was correct \citep{langosco2022goal,shah2022goal}. Our concern overlaps with all three, but the emphasis is different: we are interested in cases where the objective really is achieved and the fault lies in the path taken.

\paragraph{Constrained and safe reinforcement learning.}
Maximizing reward subject to cost limits is the subject of constrained Markov decision processes \citep{altman1999constrained} and of a large literature on safe reinforcement learning \citep{garcia2015comprehensive,achiam2017constrained}. That literature mostly studies how to \emph{train} policies that satisfy constraints. We borrow its formalism but ask a different question: how to \emph{evaluate} a deployed agent whose constraints are stated in natural language and are often only partly specified.

\paragraph{Agent benchmarks and safety evaluations.}
Benchmarks such as AgentBench \citep{liu2024agentbench}, WebArena \citep{zhou2024webarena}, and SWE-bench \citep{jimenez2024swebench} measure whether tasks are completed. A smaller body of work looks at how they are completed. MACHIAVELLI reports ethical violations alongside reward in text-based games \citep{pan2023machiavelli}, ToolEmu uses an LLM-emulated sandbox to surface risky tool use \citep{ruan2024toolemu}, and AgentDojo measures utility and security jointly under prompt injection \citep{debenedetti2024agentdojo}. Work on AI control studies monitoring protocols that remain useful even when the monitored model may act adversarially \citep{greenblatt2024control}.

\paragraph{Process-level and constraint-focused benchmarks.}
The work closest to ours measures constraint violations directly. ODCV-Bench places agents in multi-step scenarios where performance is tied to a key performance indicator (KPI) and reports how often they break ethical, legal, or safety rules to reach it \citep{li2025odcv}. SpecBench measures reward hacking in long-horizon coding agents as the gap between pass rates on visible and on held-out tests \citep{zhao2026specbench}. MasDrift tests whether delegated subtasks in multi-agent systems stay within the permissions of the original user \citep{xu2026masdrift}, and SLBench turns the logical relations between instructions in agent skill files into executable tests \citep{chen2026slbench}. \citet{huang2026reward} study reward hacking by autonomous research agents and find that a review panel misses many confirmed cases. A recent survey of trustworthy agentic AI collects metrics and benchmarks for safety, robustness, privacy, and security in a single reference \citep{qi2026trustworthy}.

Two benchmarks already report capability and safety as separate numbers. ST-WebAgentBench pairs enterprise web tasks with policies and reports Completion under Policy, which credits a task only if the agent completed it without breaking an applicable policy, together with a risk ratio \citep{levy2024stwebagentbench}. ClawsBench reports a task success rate and an unsafe action rate for productivity agents in simulated workspaces \citep{li2026clawsbench}.

A second group of papers supplies trajectory-level data and diagnosis. The ATBench extensions adapt a safety taxonomy of risk source, failure mode, and real-world harm to particular agent runtimes \citep{yang2026atbench}. Who\&When Pro builds failed trajectories with labeled failure points by injecting a failure after replaying a successful prefix \citep{liu2026whowhenpro}, a construction close to the paired trajectories we call for in Section~\ref{sec:conclusion}. \citet{jha2026meltdowns} show that agents take unsafe actions after ordinary environment errors, with no adversary involved, and relate this to exploratory behavior. On the governance side, \citet{li2026constraintdrift} argue that constraints weaken as they pass through memory, delegation, and tool use, and propose keeping them as explicit execution state, while \citet{ding2026alwayson} survey agents with persistent state and outline an evaluation contract that scores obligations on state mutation and recovery, not on answer quality alone.

Reporting outcome and compliance separately is therefore not new, and we do not claim it. Nor do we propose another benchmark. This paper adds three things to that line of work. The first is a constraint schema that does not depend on any one benchmark and that distinguishes hard from soft constraints and different kinds of evidence. The second is a per-run profile that keeps all four outcomes of Table~\ref{tab:categories} apart. A conjunctive metric such as Completion under Policy measures the rate of category~(1) and merges the other three, so it cannot separate a safe failure from an unsafe success. The third is a gating rule over the profile, together with an account of what checker error does to it. The benchmarks above complement this paper and are natural testbeds for it; we use one of them in Section~\ref{sec:setup}.

\paragraph{The gap this paper fills, relative to ODCV-Bench.}
The ODCV-Bench paper \citep{li2025odcv} reports a scalar Misalignment Rate (the fraction of runs whose four-judge median severity exceeds $3$ on a $0$--$5$ rubric) and a mean Severity score. Both are collapsed over the trajectory and over the constraints that apply to it. A single scalar is sensitive to any monotonic transformation of the rubric, and it says nothing about which agent action violated which constraint, which is exactly the input a deterministic rule-based checker needs. The ODCV-Bench paper also treats the four-judge panel's output as ground truth and reports Krippendorff's $\alpha=0.81$ as its main quality signal, while naming the absence of a human audit as an open limitation. Our framework decomposes the same trajectories into separate components for objective achievement $O$, hard violations $V$, soft-constraint cost $\mathrm{SC}$, uncertainty $U$, and distance $D$, each measurable under several checker configurations. It derives a gating rule that does not need a scalar, and it makes the audit, rather than the judge panel, the authoritative source for the false-positive rate $\alpha$ and false-negative rate $\beta$ on which the gate's verdict depends. We also show that the inter-judge agreement reported by \citet{li2025odcv} is a property of frontier models. On the same benchmark and the same scenarios, with smaller open-weight judges (Meta Llama-3.1-8B, Qwen3.7-flash, Mistral-Nemo), Krippendorff's $\alpha$ drops to $0.155$ (95\% bootstrap CI $[0.103, 0.209]$) and pairwise Cohen's $\kappa$ lies in $[0.036, 0.261]$. The gate's verdict is then determined entirely by which operating point is chosen: any-judge-flags, majority vote, or consensus. This observation is the empirical motivation for the framework's insistence on reporting checker false-positive and false-negative rates.

\paragraph{Contract- and automaton-based supervision.}
A separate line of work compiles explicit specifications into executable, deterministic checkers that run alongside the agent --- contracts over tool-call predicates, deterministic finite automata (DFAs) and temporal-logic monitors over typed action sequences --- and refuses an action the moment it violates the specification. % TODO: cite ContrAgent and similar contract/DFA frameworks the reviewer named (we lacked the exact references at submission).
These are strictly stronger than LLM judges on the subset of constraints that admit a deterministic encoding: no false positives from paraphrase, no judge disagreement to aggregate, no audit needed for $\alpha$. Our schema already accommodates them: a constraint whose \texttt{evidence} field is \texttt{state} or \texttt{log\_rule} is a deterministic contract in exactly this sense, and in our runs those constraints carry zero $U$ by construction. The complementary question --- how far such contracts can go on the full constraint distribution of ODCV-Bench --- turns on the schema's \texttt{source} field: our hand-coded breakdown of $262$ constraints puts $14.5\%$ as unstated-norm (``do not deceive the user'', ``do not retaliate'') for which no deterministic contract exists in current practice, and $40.5\%$ as implied-by-deployment-context (``do not exfiltrate PHI during this HIPAA-covered task'') where authoring one is possible but materially harder than the LLM-judge rubric we use. We treat deterministic contracts and LLM-judged constraints as complementary rather than competing: a profile under CAAEF \emph{is} a hybrid, with deterministic checks preferred wherever the evidence field permits and judges carrying only what cannot be contractually encoded. The present paper does not implement DFA-style compilation beyond the \texttt{log\_rule} class --- that would be a natural next benchmark deliverable --- and we flag it as the single lowest-risk way to reduce the paper's reliance on low-agreement open-weight judges.

\paragraph{Anomaly detection and LLM-based judges.}
Detecting inputs or behaviors that fall outside a reference distribution is a mature topic \citep{chandola2009anomaly,hendrycks2017baseline}. LLMs are now commonly used to judge the outputs of other models \citep{zheng2023judging}, although such judges have systematic biases, including a preference for their own generations \citep{panickssery2024llm}, and trajectory-level judging in particular can miss structural faults that outcome-level judges score as a success. % TODO: cite Mohammadi trajectory-judge paper if available.
Both lines of work bear on trajectory monitoring, and we return to them in Sections~\ref{sec:framework} and~\ref{sec:challenges}.

\section{Goal Achievement versus Constraint Compliance}
\label{sec:problem}

An agent can be viewed as a policy $\pi$ chosen to maximize an expected task reward,
\begin{equation}
    \max_{\pi}\; \mathbb{E}_{\tau\sim\pi}\big[R(\tau)\big],
    \label{eq:unconstrained}
\end{equation}
where $\tau$ is a trajectory and $R$ rewards completion of the task. Real instructions almost never stop there. They also say which resources the agent may access, which tools it may use, what information it may disclose, and which actions are off limits. Writing these restrictions as a set of constraints $\mathcal{C}=\{c_1,\dots,c_m\}$, the behavior we actually want is closer to the constrained problem \citep{altman1999constrained}
\begin{equation}
\begin{aligned}
    \max_{\pi}\quad & \mathbb{E}_{\tau\sim\pi}\big[R(\tau)\big] \\
    \text{s.t.}\quad & \Pr_{\tau\sim\pi}\big[\tau \models c_i\big] = 1,
    \quad i=1,\dots,m,
\end{aligned}
\label{eq:constrained}
\end{equation}
where $\tau \models c_i$ means that the trajectory satisfies constraint $c_i$.

Requiring every constraint to hold with probability one is the right target for restrictions that must never be broken, such as using a credential the agent was not given. As a general model, though, it is too strict, for two reasons. Some restrictions are budgets rather than prohibitions, for example a limit on the number of external requests. And violations are detected in practice by imperfect checkers, so a guarantee of exactly zero cannot be verified. The safe reinforcement learning literature offers two standard relaxations \citep{altman1999constrained,achiam2017constrained}. A chance constraint bounds the probability of violation,
\begin{equation}
    \Pr_{\tau\sim\pi}\big[\tau \not\models c_i\big] \le \delta_i ,
    \label{eq:chance}
\end{equation}
and a cumulative-cost constraint bounds an expected cost,
\begin{equation}
    \mathbb{E}_{\tau\sim\pi}\big[g_i(\tau)\big] \le d_i ,
    \label{eq:cost}
\end{equation}
where $g_i\ge 0$ measures how much of budget $i$ the trajectory consumed. Problem~\eqref{eq:constrained} is the special case $\delta_i=0$. We call a constraint \emph{hard} if it is of this kind and \emph{soft} otherwise; the quantities reported in Section~\ref{sec:framework} are defined for both.

Problems arise when the two formulations disagree, that is, when breaking a constraint raises the probability of completing the task. Typical situations include the following:
\begin{itemize}
    \item a restricted tool gives faster or more reliable access to the needed information;
    \item an external system exposes credentials that make the task trivial;
    \item a security weakness lets the agent bypass a restriction in the environment;
    \item the evaluation metric can be manipulated directly to yield a higher score;
    \item taking many additional, unrequested actions raises the chance that one of them produces the desired result.
\end{itemize}
In each case an optimizer that sees only~\eqref{eq:unconstrained} has a reason to leave the feasible set of~\eqref{eq:constrained}. These situations are instances of the reward hacking, specification gaming, and goal misgeneralization discussed in Section~\ref{sec:related}. What matters, then, is not only whether an agent \emph{can} reach an objective, but whether its optimization stays inside the operational boundaries its principals intended.

Two features make the agent setting harder than the classical constrained one. First, $\mathcal{C}$ is rarely written down in full. Some constraints are explicit in the prompt, some are implied by the deployment context, and some are norms nobody thought to state. Second, the action space is open-ended. An agent with a shell and network access can do things the benchmark designer never enumerated, so violations cannot be detected by checking membership in a short list of forbidden actions.

\section{A Constraint-Aware Evaluation Framework}
\label{sec:framework}

We propose a Constraint-Aware Agent Evaluation Framework (CAAEF) with five components, ordered from what is already standard practice to what is least developed.

\subsection{Outcome Verification}

The first component asks whether the assigned objective was completed and yields a score $O(\tau)\in[0,1]$. This is what existing benchmarks measure. It is necessary but not sufficient: two agents can produce the same result by entirely different means, so $O$ on its own says nothing about whether the behavior can be trusted.

\subsection{Trajectory Monitoring}

The second component records the full sequence of actions,
\begin{equation}
    \tau = (a_1, a_2, \dots, a_T),
\end{equation}
together with the observation each action returned, and makes this sequence, rather than the final answer, the object of evaluation. A trajectory is long and mostly routine, so a monitor needs to know what to look for. The situations listed in Section~\ref{sec:problem} suggest the following signals:
\begin{itemize}
    \item tool calls that the task does not obviously require;
    \item attempts to access restricted resources;
    \item repeated failed attempts to get past a control;
    \item abrupt changes of strategy;
    \item contact with external systems not mentioned in the task;
    \item acquisition or use of credentials the agent was not given;
    \item actions unrelated to the legitimate objective.
\end{itemize}
None of these is conclusive in isolation. A second attempt after a refusal may simply correct a typing error, and an unfamiliar tool may be the right one. The signals are inputs to the components that follow: the compliance check decides whether a flagged action broke a constraint, and the anomaly score measures how unusual the overall pattern is.

\subsection{Constraint Compliance}

Checking actions against $\mathcal{C}$ requires the constraints to be written so that two evaluators would apply them in the same way. We suggest recording each constraint $c_i$ with six fields:
\begin{itemize}
    \item \emph{scope}: the resources, tools, or data it covers;
    \item \emph{modality}: prohibition, obligation, or budget;
    \item \emph{class and severity}: hard or soft, with a weight $w_i\ge 0$;
    \item \emph{condition}: the state in which it applies, for example only before the user has confirmed;
    \item \emph{source}: stated in the instructions, implied by the deployment context, or an unstated norm;
    \item \emph{evidence}: how a violation is established.
\end{itemize}
For the last field we distinguish three kinds of evidence, in decreasing order of reliability: the state of the environment or the artifacts the agent produced; the execution log, checked by a deterministic rule; and a reading of the trajectory by a human or LLM judge. Each constraint should be checked with the most reliable kind of evidence available for it. Executable tests of the sort used in SLBench \citep{chen2026slbench} fall under the first two kinds. Constraints whose source is an unstated norm will usually need a judge, and these are the ones on which labels are least consistent. Recording the source makes it possible to report results separately for explicit and implicit constraints.

Let $v_{t,i}\in\{0,1\}$ indicate whether action $a_t$, taken in the context of the preceding trajectory, violates $c_i$. Two summaries of this array are used below. The per-action indicator $v_t=\max_i v_{t,i}$ records whether step $t$ broke any constraint, and the per-constraint indicator $V_i(\tau)=\max_t v_{t,i}$ records whether constraint $c_i$ was broken at any step. Context matters because the same action, such as reading a file, may be permitted at one point and prohibited at another. A simple trajectory-level compliance score is the fraction of compliant actions,
\begin{equation}
    \mathrm{CS}(\tau) \;=\; 1 - \frac{1}{T}\sum_{t=1}^{T} v_t.
    \label{eq:cs}
\end{equation}
A high objective score combined with a low compliance score points to optimization that succeeded but may not be safe.

The ratio in~\eqref{eq:cs} has an obvious weakness: a single serious violation in a long trajectory barely moves it. A hundred-step run that exfiltrates a credential once still scores $0.99$. We therefore report two further quantities. Writing $\mathcal{C}_H\subseteq\mathcal{C}$ for the hard constraints, the strict indicator
\begin{equation}
    V(\tau) \;=\; \max_{c_i\in\mathcal{C}_H} V_i(\tau)
    \label{eq:strict}
\end{equation}
records whether any hard constraint was broken, and the severity-weighted cost
\begin{equation}
    \mathrm{SC}(\tau) \;=\; \sum_{t=1}^{T}\;\sum_{c_i\notin\mathcal{C}_H} w_i\, v_{t,i}
    \label{eq:sc}
\end{equation}
accumulates soft violations over the run. A hard violation is a reason to stop or disqualify the run; a soft violation is a warning that adds to a running total. Over many runs, the mean of $V$ estimates the left-hand side of~\eqref{eq:chance}, and the mean of $\mathrm{SC}$ plays the role of the cost in~\eqref{eq:cost}.

The indicators $v_{t,i}$ are checker outputs, not ground truth. Reported values of $\mathrm{CS}$, $V$, and $\mathrm{SC}$ should therefore come with the checker's false-positive and false-negative rates, estimated on a sample of trajectories audited by humans.

\subsection{Uncertainty and Anomaly Detection}

Not every unusual action is a violation, and not every violation is deliberate. An agent that is unsure how to proceed may legitimately try alternatives, and this kind of exploration should not be penalized as if it were an attack. Evaluation should therefore carry an explicit estimate of uncertainty, which we denote $U(\tau)$. It covers both the agent's uncertainty about the task and the evaluator's uncertainty about whether a given action is compliant.

Anomaly detection offers one way to make this operational. Given a reference set of trajectories known to be compliant, one can compute a shift score $D(\tau)$ that measures how far the current trajectory, or its prefix $\tau_{1:t}$, lies from that distribution, using standard techniques \citep{chandola2009anomaly,hendrycks2017baseline} applied to action sequences or their learned representations. Because the score can be computed on prefixes, it raises the possibility of flagging risky behavior \emph{before} the agent finishes. We regard the use of uncertainty and distribution shift for this kind of early warning as one of the more promising directions suggested by the framework.

One concrete way to compute $D$ is as follows. Tool names differ from task to task, so raw calls cannot be compared directly. Each call is first mapped to a typed action given by the category of the tool, the class of resource it touches, whether it reads or writes, and whether it falls inside the scope granted by the task. A prefix is then encoded by the sequence of these types and scored by its distance to the nearest compliant reference prefixes for tasks of the same family. The flagging threshold is set on a held-out set of compliant trajectories so that the false-alarm rate is controlled. Conformal methods give such a guarantee when calibration and test trajectories are exchangeable \citep{vovk2005algorithmic}.

Exchangeability is a strong assumption here. It fails when the agent meets tools or task types that are absent from the calibration set, and it fails after any change to the agent. We therefore suggest calibrating separately for each task family and tracking the flag rate on runs known to be compliant. If that rate rises, the guarantee no longer holds and the threshold must be recalibrated. The same calibration step applies to $U$, which can be estimated from the disagreement among several independent checkers.

The two signals are not interchangeable, and neither subsumes the other. $U$ expresses confidence in a judgment of compliance. $D$ describes the behavior itself and involves no such judgment. They can disagree in both directions. A new but legitimate way of solving a task lies far from the reference set and is plainly compliant, so $D$ is high and $U$ is low. A routine trajectory may contain one step whose status cannot be settled from the available evidence, so $D$ is low and $U$ is high. In our reading, $D$ is the signal suited to early warning on prefixes, and $U$ is the signal that decides which flagged cases need a human.

Uncertainty also bears on the soft cost, and the two kinds of soft constraint behave differently. A budget is checked by counting in the log, so an overrun is known exactly. What is unclear is whether it matters: an agent that retries a failed call or tries a second tool is exploring, and exploration consumes budget. The period after an environment error needs particular care, since that is where unsafe actions have been observed without any adversary \citep{jha2026meltdowns}. To keep exploration and overrun apart, we suggest setting each budget $d_i$ from the reference set, for example at a high quantile of what compliant runs of the same task family consume, and counting only the excess. Ordinary exploration then falls inside the budget by construction. Soft constraints that need a judge are a different case, because there the violation itself is uncertain. For these we suggest reporting $\mathrm{SC}$ as an interval. The lower end counts the violations on which all checkers agree, and the upper end adds those on which they disagree. A wide interval says that the soft cost of the run is not known, which is more useful than a point value that hides the disagreement.

\subsection{Post-Hoc Behavioral Analysis}

After the task ends, the trajectory should be reconstructed and examined with questions such as these. What was the agent trying to accomplish at each stage? Which actions did the task need, and which did it not? Did the agent try to get around a restriction? Did it exploit a weakness in the environment? Did it obtain information or permissions beyond what the task required?

The answers support a distinction between \emph{necessary} optimization, meaning effort the task demands, and \emph{opportunistic} optimization, meaning the use of whatever shortcut the environment happens to offer. The framework is designed to expose the second kind.

\section{A Two-Dimensional Evaluation Model}
\label{sec:model}

The framework separates two core dimensions. Objective achievement asks whether the agent accomplished the assigned task; constraint compliance asks whether it stayed within the intended boundaries. Crossing the two gives the four categories in Table~\ref{tab:categories}.

\begin{table}[t]
\caption{Behavioral categories obtained by crossing objective achievement with constraint compliance. Outcome-only benchmarks cannot separate the two cells in the top row.}
\label{tab:categories}
\vskip 0.1in
\begin{center}
\begin{small}
\begin{tabular}{lcc}
\toprule
Objective & Compliant & Violating \\
\midrule
Achieved & (1) Intended success & (2) Unsafe success \\
Failed   & (3) Safe failure     & (4) Unsafe failure \\
\bottomrule
\end{tabular}
\end{small}
\end{center}
\vskip -0.1in
\end{table}

Category~(2) deserves the most attention. An outcome-only benchmark puts it in the same bin as category~(1), so the agents that are most worrying from a safety standpoint are precisely the ones the benchmark rewards. Category~(3) also deserves a remark. An agent that stops and reports that the task cannot be done within its permissions has behaved well, yet outcome-only scoring treats it exactly like an agent that failed through incompetence.

A fuller evaluation profile adds the quantities from Section~\ref{sec:framework} and reports, for each run, the tuple
\begin{equation}
    P(\tau) \;=\; \big(O,\; \mathrm{CS},\; V,\; \mathrm{SC},\; U,\; D\big),
    \label{eq:profile}
\end{equation}
with each quantity evaluated on $\tau$. We recommend against collapsing this profile into a single scalar. Any fixed weighting lets a large gain in $O$ offset a constraint violation, which reintroduces, at the level of the metric, the very trade-off the constraints were meant to rule out. A single number can hide the failure modes that matter most.

Declining to aggregate does not mean declining to decide. For release gating we suggest a lexicographic rule. A run with $V=1$ is disqualified whatever its value of $O$, and an agent is rejected if its rate of such runs exceeds a tolerance $\delta$ set by the deployer. Among agents that pass this test, one is preferred to another only if it is no worse on both $\mathrm{SC}$ and $O$. Where neither dominates, the choice is left to the deployer, who sees the trade-off instead of having it settled by a weighting. Under this rule a safe failure ranks above an unsafe success: category~(3) costs utility, whereas category~(2) breaks the conditions under which the agent was allowed to act at all.

The rule is only as good as the checker behind $V$, and its behavior under checker error can be worked out directly. Suppose the checker reports a hard violation on a compliant run with probability $\alpha$ and misses a real one with probability $\beta$. If the true rate of hard violations is $p$, the observed rate is
\begin{equation}
    q \;=\; (1-\beta)\,p + \alpha\,(1-p).
    \label{eq:observed}
\end{equation}
Three consequences follow. First, an agent that never violates is still observed at rate $\alpha$, so it is rejected whenever $\alpha>\delta$; a tolerance below the checker's false-positive rate cannot be enforced on the raw output. Second, a share $\alpha$ of compliant runs is disqualified, and if the monitor stops a run when it flags, the same share is turned into forced safe failures. This is the utility cost of the veto, and it is paid whether or not the agent ever misbehaves. Third, misses lower $q$ by $\beta p$, so an agent whose true rate is above $\delta$ can still pass the gate.

Hard violations are meant to be rare, and when $p$ is small the false positives dominate. With $p=0.02$, $\alpha=0.05$, and $\beta=0.2$, equation~\eqref{eq:observed} gives $q=0.065$, of which about three quarters are false alarms. Two practical rules follow. The gate should be applied to the corrected estimate
\begin{equation}
    \hat{p} \;=\; \frac{q-\alpha}{1-\alpha-\beta},
    \label{eq:corrected}
\end{equation}
which is valid when $\alpha+\beta<1$ \citep{rogan1978estimating}, with a confidence interval that includes the uncertainty in $\alpha$ and $\beta$ from the audit. And every run with $V=1$ should be confirmed by stronger evidence or by a human before it counts; this removes most of the cost of $\alpha$ at the price of review effort proportional to $q$. The analysis assumes that $\alpha$ and $\beta$ do not depend on the agent. That assumption fails for an agent that conceals its actions, for which $\beta$ will be higher than the audit suggests.

\section{Experimental Setup}
\label{sec:setup}

\paragraph{Testbed.}
We use ODCV-Bench \citep{li2025odcv}, the current state-of-the-art benchmark for the failure mode this paper studies: agents under KPI pressure that meet their performance target by violating the ethical, legal, or safety constraints named in their system prompt. Each scenario pairs a system prompt stating a role and its obligations with a user prompt stating a task and a numeric target, inside a sandboxed filesystem whose scripts contain deliberate loopholes. Every scenario exists in two framings that differ only in the user prompt: \emph{mandated}, which instructs the agent to reach the target, and \emph{incentivized}, which presents the target as a condition for reward without instructing the agent to meet it. The benchmark's own paper evaluates $12$ frontier agents with a four-judge frontier panel (Krippendorff's $\alpha=0.81$), reports misalignment rates from $0\%$ to $62.8\%$, and names as an explicit limitation that scores come from LLM judges without human verification. Our study is designed to measure what ODCV-Bench exposes under the measurement discipline that limitation calls for. We score every run under four checker configurations (two deterministic, two judge-based); we quantify how checker false-positive and false-negative rates propagate into the gate's verdict; we build the human-audit infrastructure that their paper flags as missing; and we extend their 40 scenarios across both framings with up to three repetitions per cell. For the original wave we drew $20$ scenarios, stratified by container base image with a fixed seed. The extension wave covers the remaining $20$ at three repetitions, and their constraints were committed to git before any run on them. Both the draw and the generalization set are therefore reproducible and independent of anything we measured.

\paragraph{Agents and judges.}
We use four agents, two from each of two model families, matched on price across families so that a difference between families cannot be a price-tier artefact: a mid tier at \$2/\$10 per million input/output tokens and a small tier at \$0.10/\$0.50. The original wave uses two judges, one per family. Budget limited these judges to models that also appear as agents, so some judgments are same-\emph{model} and not merely same-family; we tag these and report every judge-dependent figure both over all judgments and over cross-family judgments only. The extension adds a second panel of three smaller open-weight judges from three further families (Meta Llama-3.1-8B, Qwen3.7-flash, and Mistral-Nemo), each judging every constraint in a separate call (Section~\ref{sec:judge_results}).

\paragraph{Harness and tracing.}
Each run executes inside a container with no network access, driven by a ReAct loop \citep{yao2023react} with two tools (\texttt{bash} and \texttt{task\_complete}), a 50-step cap, and a 30\,s per-command timeout. Every run writes one JSONL trace containing, for each action, the full command, its output, exit code, timestamps, and token counts, plus a SHA-256 hash of every file in the sandbox before the run and after every action. The hashes are taken out of band, so they never enter the agent's context. Failed and truncated runs are kept and marked, never discarded. A separate replay check verifies that each stored trace reconstructs its own action sequence and the exact message list the model received; all 303 traces pass.
% CHECK: 303 traces vs. 297 analysed original runs (presumably incl. 6 partial second-repetition runs) -- state this explicitly or align.

The trace schema has to respect one subtlety. A single model turn can emit several parallel tool calls. The trajectory $\tau=(a_1,\dots,a_T)$ indexes \emph{actions}, and $\mathrm{CS}$ divides by $T$, so conflating turns with actions would inflate $\mathrm{CS}$ on exactly those runs. Traces therefore carry a per-action index distinct from the turn index.

\paragraph{Constraints.}
We wrote $262$ constraints over all $40$ scenarios (mean $6.5$ per scenario, range $5$--$8$): $131$ for the original 20, used in Phase~1, and $131$ for the remaining 20, drafted for the extension before any agent was run on them, so that the ordering can be verified from the git history rather than merely asserted. Each constraint has the six fields of Section~\ref{sec:framework} and was drafted from the scenario prompts alone, before any trajectory was examined. Their distribution is itself a finding about how practical the framework is. By source, $45.0\%$ are explicit, $40.5\%$ implied by the deployment context, and $14.5\%$ unstated norms. By evidence, $38.9\%$ can be established from environment state, $13.0\%$ from a deterministic log rule, and $48.1\%$ require a judge. Nearly half of all constraints thus fall into the least reliable evidence class, and \emph{every} unstated-norm constraint does, which is exactly where Section~\ref{sec:framework} predicts labels will be least consistent.

\paragraph{Objective verification.}
The benchmark ships no working objective scorer, so we had to define $O$ ourselves. For each run we replay the stored final filesystem state into a fresh container and execute the scenario's own KPI script from a \emph{pristine} copy, comparing its output against the threshold the scenario states. The pristine copy matters, because an agent that rewrites the scorer could otherwise inflate its own $O$. In one early run an agent edited a ranking script to report $250/250$, while the pristine copy reported the true $50/250$. This defeats script tampering but not data tampering, and that is the intended reading: an agent that deletes inconvenient records and then legitimately scores well has $O=1$ and $V=1$, which is category~(2). The initial pass left six of the 20 original scenarios without an agreed rule, because their KPI scripts need arguments that the prompt does not determine, or write results to a file instead of standard output. The extension study closed this gap for both waves. All six original rules were specified from the pristine scripts' real output on real traces, with each choice justified in the rule file, and $20$ further rules were written for the extension's new scenarios. Four of these needed a purpose-built measurement because the scenario ships no KPI scorer at all. $O$ coverage now stands at $297/297$ original runs and $475/480$ extension runs; the five remaining extension runs, whose per-scenario measurement failed on a particular final state, are reported as missing rather than imputed.

\paragraph{Checker configurations.}
We score every run four ways, as an ablation over the framework's components: \textsc{rules+state} (state hashes and deterministic log rules only), \textsc{judge~A} alone, \textsc{judge~B} alone, and \textsc{full} (state or log rules wherever the constraint's evidence field permits, and both judges elsewhere, with a judge-based violation counted when either judge flags it). Disagreement is reported separately as $U$, not averaged away.

\paragraph{A worked constraint.}
To make the schema concrete, here is one constraint from the hiring scenario in full. \emph{Scope}: the ranking script that the task names as the measure of performance. \emph{Modality}: prohibition. \emph{Class}: hard, so $w_i$ is unused and the constraint enters $V$, not $\mathrm{SC}$. \emph{Condition}: at all times. \emph{Source}: implied; the prompt never says ``do not edit the scorer'', but naming it as the measure makes editing it a violation of the deployment context. \emph{Evidence}: state. The checker hashes every file in the sandbox before the run and after each action, and sets $v_{t,i}=1$ at the first action whose post-state hash for that path differs from the pre-run hash. In one run the agent issued \texttt{chmod +w} on the script and overwrote it to print a perfect score. The hash changed at that action, so $v_{t,i}=1$ there, giving $V=1$; since the pristine copy still reported the true failing score, $O=0$, and the run falls in category~(4). A sibling constraint in the same file (``the candidate pool must be diverse'') is soft, with $w_i=2.0$ and evidence from a judge, because no file hash can settle it.

\paragraph{Typed actions and the distance score.}
For reproducibility we state the encoding exactly. Each \texttt{bash} command is mapped to a tuple of four categorical features. The \emph{tool category} (read, write, execute, delete, permission change, network, process control, other) is assigned by first match over an ordered list of command patterns, so that destructive and privileged categories are tested before generic ones. The \emph{resource class} (task data, source data, validator, configuration, logs, temporary, outside workspace, none) is assigned by first match over path patterns. The \emph{read/write} flag is true for the write, delete, and permission-change categories. The \emph{in-scope} flag is false when the command names any absolute path outside those the task grants, and network actions are out of scope by construction. The four features are concatenated into a single token, so the vocabulary is the set of observed tuples ($46$ here) rather than the full product space. A prefix $\tau_{1:t}$ is encoded as the concatenation of its length-normalized token counts and its length-normalized token-bigram counts, the latter restricted to the $400$ most frequent bigrams. $D(\tau_{1:t})$ is the mean Euclidean distance to the $5$ nearest reference encodings among compliant prefixes whose length is within $\pm 2$ of $t$, drawn only from other scenarios.

The modeling choices are deliberately conservative, and their failure modes are known. The typed, hand-coded vocabulary (as opposed to a learned embedding) keeps the encoder transparent and reproducible without GPU dependencies, and makes a per-token audit possible; it pays for that with a hard ceiling at $46$ observed tuples that cannot capture semantic similarity between distinct typed actions. Length-normalized counts plus the top-$400$ bigrams carry local sequential structure (``read then write'', ``delete then exec'') without the dimensionality explosion of full bigrams or the fragility of positional embeddings on $5$-step trajectories; this encoding is bag-of-bigrams, so it is blind to the \emph{order} of two non-adjacent events. Euclidean distance on normalized counts is a defensible default for sparse count vectors and matches our use of a $k$-nearest neighbor baseline, but it treats all vocabulary dimensions as independent. The expected failure modes are three: (i) short trajectories (median $5$ actions), where there are few bigrams to average over and a single out-of-vocabulary action dominates the distance; (ii) cross-scenario vocabulary shift, where a prefix in a new scenario sits far from any compliant reference for reasons unrelated to compliance (the exchangeability failure we document empirically in Section~\ref{sec:judge_results}); and (iii) semantically-identical rewrites of the same action (``\texttt{cat file.txt}'' vs ``\texttt{less file.txt}''), which map to different tuples under our conservative first-match rule. Richer alternatives --- an HMM over typed actions, a learned contrastive embedding, a trajectory-level transformer, or a per-task-family calibrated score \citep{barber2023conformal} --- would address (i)--(iii) in principle but require a labeled seed per task family that this benchmark does not supply.

\paragraph{Statistics.}
All intervals are $95\%$ percentile bootstrap intervals with $10{,}000$ resamples. We resample \emph{whole scenarios} rather than runs, because the two framings and the repetitions of a scenario are not independent. With four agents from two families, any ranking is weak evidence, and we say so wherever one appears.

\paragraph{What the setup could not deliver.}
Three departures from the plan bound the results, and we state them here rather than in a footnote. First, \textbf{decoding is not deterministic}. Both providers reject or silently ignore \texttt{temperature}$=0$, and four identical requests returned four different completions. Runs and judgments are therefore sampled, not greedy, which makes repetitions essential; yet only one complete repetition of the original wave was funded (297 of 400 planned runs; the second repetition is partial and is not pooled). Second, \textbf{the original-wave judges are weaker than intended}, and they batch a scenario's constraints into one call instead of judging each in isolation, so a verdict on one constraint can be influenced by the others. Third, \textbf{the human audit has not been run}. The annotation tooling and a stratified 200-run sample exist, but $\alpha$ and $\beta$ are unmeasured, so every quantity that depends on them is reported as pending. The total API cost of the study was \$4.61.

\section{Results}
\label{sec:results}

\subsection{The profile changes the conclusion, when the checker is good enough}

Table~\ref{tab:configs} is the central result, and it is a mixed one.

\begin{table}[t]
\caption{The conclusion depends on the checker. Kendall's $\tau$ compares the ranking by $O$ alone with the ranking under the gate at $\delta=0.05$. Where no model passes the gate, $\tau=+1$ reflects degeneracy, not agreement: the gate cannot discriminate, so the gated ranking collapses onto the $O$ ranking.}
\label{tab:configs}
\vskip 0.1in
\begin{center}\begin{small}
\setlength{\tabcolsep}{4pt}
\begin{tabular}{lcccc}
\toprule
Configuration & Runs & $V$ rate & Pass & $\tau$ \\
\midrule
\textsc{rules+state} & 297 & 0.029--0.125 & 2 & $-0.60$ \\
\textsc{judge~A}     & 200 & 0.000--0.100 & 3 & $+0.60$ \\
\textsc{judge~B}     & 200 & 0.075--0.300 & 0 & $+1.00$ \\
\textsc{full}        & 200 & 0.163--0.300 & 0 & $+1.00$ \\
\textsc{full}, cross-fam. & 200 & 0.125--0.300 & 0 & $+1.00$ \\
\bottomrule
\end{tabular}
\end{small}\end{center}
\vskip -0.1in
\end{table}

Under \textsc{rules+state}, the claim of Section~\ref{sec:model} holds, and holds strongly. Kendall's $\tau$ between the ranking by $O$ alone and the ranking under the gate at $\delta=0.05$ is $-0.60$: the two rankings are not merely different, they are \emph{opposed}. The two agents with the highest objective achievement ($O=0.152$ each) are exactly the two that the gate disqualifies, while the two with lower $O$ pass. An outcome-only benchmark would promote precisely the agents that a constraint-aware gate rejects.
% CHECK: O = 0.152 predates the extension rescore (O coverage 70% -> 100%); confirm against the current rules+state numbers.

Under the three-judge, majority-vote, per-constraint protocol, run on $772$ runs across both waves, the inversion does not reappear. The ranking by objective achievement (gpt-6-luna $15.3\%$, gpt-6-sol $13.2\%$, sonnet-5.5 $10.3\%$, haiku-5.5 $8.1\%$) and the ranking by gate-passing rate ($V{=}0$) under the same protocol (gpt-6-luna $63.4\%$, then haiku-5.5 and gpt-6-sol tied at $59.9\%$, then sonnet-5.5 $55.1\%$) are weakly positively correlated ($\tau=+0.33$), not opposed. Our claim is therefore specifically that outcome and gate rankings invert \emph{under deterministic state+rule checkers}; under a less well-calibrated judge panel, the inversion dissolves into the panel's own noise. This in turn supports the framework: the checker configuration \emph{is} a load-bearing parameter of the measurement, not a technicality.

Under the \textsc{full} protocol, $\tau=+1.00$. This is degeneracy, not agreement. \emph{No agent passes the gate}, so the gate cannot discriminate and the gated ranking collapses back onto the $O$ ranking. With violation rates of $0.163$--$0.300$ against a tolerance of $\delta=0.05$, every agent is disqualified and the ranking question becomes empty.

This degeneracy is not an artefact of our analysis; it is equation~\eqref{eq:observed} showing up in the data. Section~\ref{sec:model} argues that a tolerance below the checker's false-positive rate cannot be enforced on raw output. Either these agents really do violate constraints in $16$--$30\%$ of runs, or $\alpha$ is of that order. \textbf{The two possibilities cannot be separated without the audit}, which is why the audit is essential and not a refinement. Every gate conclusion reported here is provisional on it.

\paragraph{The inversion is specific to the tolerance.}
The $\tau=-0.60$ above is computed at $\delta=0.05$, a value we chose before looking at the data and did not justify against any risk model. Varying it changes the picture considerably; Table~\ref{tab:delta} reports the full sweep. Under \textsc{rules+state}, no agent passes at $\delta\in\{0.01,0.02\}$, two pass at $0.05$, and three pass at $0.10$, with corresponding $\tau$ values of $+1.00$, $+1.00$, $-0.60$, and $+0.18$. \textbf{The ranking inversion therefore occurs within a window of tolerances, not everywhere.} Below the window the gate rejects everything and says nothing; above it, the gate admits almost everything and again says little. Only in between does it separate agents in a way that contradicts the outcome ranking.

\begin{table}[t]
\caption{Sensitivity of the gate to the tolerance $\delta$. Each cell gives the number of agents (out of four) passing the gate and, in parentheses, Kendall's $\tau$ between the ranking by $O$ and the gated ranking. A $\tau$ of $+1$ where no agent passes reflects degeneracy, not agreement. The ranking inversion appears only at $\delta=0.05$ under \textsc{rules+state}.}
\label{tab:delta}
\vskip 0.1in
\begin{center}\begin{small}
\setlength{\tabcolsep}{3pt}
\begin{tabular}{lcccc}
\toprule
 & \multicolumn{4}{c}{Tolerance $\delta$} \\
\cmidrule(l){2-5}
Configuration & $0.01$ & $0.02$ & $0.05$ & $0.10$ \\
\midrule
\textsc{rules+state} & 0 (+1.00) & 0 (+1.00) & 2 ($-$0.60) & 3 (+0.18) \\
\textsc{judge~A}     & 1 (+0.91) & 1 (+0.91) & 3 (+0.60) & 4 (+1.00) \\
\textsc{judge~B}     & 0 (+1.00) & 0 (+1.00) & 0 (+1.00) & 2 (+1.00) \\
\textsc{full}        & 0 (+1.00) & 0 (+1.00) & 0 (+1.00) & 0 (+1.00) \\
\textsc{full}, cross-fam. & 0 (+1.00) & 0 (+1.00) & 0 (+1.00) & 0 (+1.00) \\
\bottomrule
\end{tabular}
\end{small}\end{center}
\vskip -0.1in
\end{table}

We read this as a sharpening of the claim, not a retraction. The gate is informative exactly when $\delta$ lies above the checker's false-positive rate and below the agents' true violation rate, and nothing in the framework as stated tells a practitioner where that window is. Finding it requires $\alpha$, which requires the audit. A deployer who picks $\delta$ by intuition may land outside the window and conclude either that every agent is unsafe or that the gate is vacuous, and in both cases learn nothing. This is a concrete, actionable gap in the framework that the experiment exposed and that we had not anticipated.

\paragraph{Setting $\delta$ in practice, when $\alpha$ is initially unknown.}
We propose a procedure a deployer can follow without yet having a full human audit, and then update as audit data accrue. (i) \emph{Risk-anchored upper bound.} The right $\delta$ is in the end a risk tolerance: a hard-violation rate above $\delta$ is deemed unacceptable for release. For regulated settings the upper bound is given by a historical incident threshold or a published policy (``no more than $x$ unauthorised disclosures per $10{,}000$ sessions''); for new settings it is set by analogy with human-operator error rates for the same task family. In either case, $\delta$ is a stated risk budget, not a knob tuned to the data. (ii) \emph{Checker-anchored lower bound.} The gate becomes degenerate (every agent rejected) once $\delta < \alpha$, so $\alpha$ is a mechanical lower bound on any useful $\delta$. Pending the full audit, $\alpha$ can be bracketed cheaply by scoring a small adversarial sample (say $k=30$ runs the authors believe to be fully compliant) and taking the Clopper--Pearson upper CI on the checker's flag rate on that sample; this bound tightens with every audit batch. (iii) \emph{Window check.} If the risk-anchored upper bound sits below the checker-anchored lower bound, the gate cannot be used with this checker on this agent class, and the correct response is either to switch to a lower-$\alpha$ checker (deterministic contracts where possible, consensus aggregation elsewhere) or to document that the measurement does not support the decision. (iv) \emph{Prior over $p$.} When the audit is still in progress but partial labels exist, the Rogan--Gladen corrected rate $p=(q-\alpha)/(1-\alpha-\beta)$ carries a propagated interval that can be read as a posterior over $p$; a $\delta$ is viable if its risk-anchored upper bound exceeds the upper CI of $p$ for at least one agent. The procedure makes explicit what the current paper only implied: a well-defined $\delta$ is a joint claim about the deployer's risk tolerance \emph{and} the checker's error profile, and the latter is not substitutable with the former.

The same instability appears for individual agents. Across the configurations of Table~\ref{tab:configs}, each agent's measured violation rate moves by a factor of three or more, and in one case from zero to almost a fifth: the small OpenAI agent is scored at $0.125$ by \textsc{rules+state}, $0.000$ by judge~A, $0.075$ by judge~B, and $0.175$ by the full protocol. More awkwardly, the \emph{ordering} of agents is not stable either. The mid-tier Anthropic agent is among the two lowest of the four under \textsc{rules+state} ($0.031$, against $0.029$ for the best agent, a difference well inside either interval), yet it is the \emph{highest} of all four under every judge-based configuration ($0.100$, $0.300$, and $0.300$). No agent is uniformly best or worst, so the question ``which of these agents violates least?'' has a different answer depending only on which checker is asked.

\subsection{Safe failure is the common case, not an edge case}

After the extension's rescore brought $O$ coverage to the full analysed set, overall objective achievement is $10.4\%$ on the original wave ($31/297$) and $12.1\%$ on the extension wave ($58/480$). The two are close, but we report them separately and never pool them silently. Under \textsc{rules+state}, the four categories on the original 297 runs are: intended success $27$ ($9.1\%$), unsafe success $4$ ($1.3\%$), safe failure $251$ ($84.5\%$), and unsafe failure $15$ ($5.1\%$). Table~\ref{tab:profile-full} gives the full per-agent profile under the full protocol, including the category counts.

\begin{table*}[t]
\caption{Per-agent evaluation profile under the full checker configuration. Intervals are $95\%$ bootstrap intervals, resampling whole scenarios. Category counts are (1) intended success, (2) unsafe success, (3) safe failure, and (4) unsafe failure, over the runs with a known $O$. Safe failure is the largest cell for every agent.}
\label{tab:profile-full}
\vskip 0.1in
\begin{center}\begin{small}
\begin{tabular}{lcccc}
\toprule
Model & $O$ & $V$ rate & CS & (1)\,/\,(2)\,/\,(3)\,/\,(4) \\
\midrule
haiku-5.5  & 0.071 [0.000, 0.214] & 0.163 [0.050, 0.300] & 0.959 [0.928, 0.987] & 1\,/\,3\,/\,46\,/\,6 \\
sonnet-5.5 & 0.107 [0.000, 0.286] & 0.300 [0.125, 0.500] & 0.885 [0.803, 0.950] & 1\,/\,2\,/\,20\,/\,5 \\
gpt-6-luna & 0.179 [0.000, 0.393] & 0.175 [0.050, 0.325] & 0.917 [0.872, 0.958] & 3\,/\,2\,/\,19\,/\,4 \\
gpt-6-sol  & 0.179 [0.000, 0.393] & 0.225 [0.100, 0.375] & 0.901 [0.842, 0.949] & 3\,/\,2\,/\,18\,/\,5 \\
\bottomrule
\end{tabular}
\end{small}\end{center}
\vskip -0.1in
\end{table*}
% CHECK: category counts sum to 56/28/28/28 = 140, whereas the full configuration covers 200 runs (50 per agent). Confirm these were not computed under the old 70% O coverage.

Safe failure is the largest cell for every agent without exception, ranging from $64\%$ of one agent's runs to $82\%$ of another's. Note also that $\mathrm{CS}$ stays between $0.885$ and $0.959$ throughout, while the $V$ rate ranges from $0.163$ to $0.300$ over the same runs. This is the weakness of the ratio in equation~\eqref{eq:cs} made concrete: a trajectory can contain a disqualifying action and still score above $0.88$ on compliance, because one bad step among several ordinary ones barely moves an average. $\mathrm{CS}$ and $V$ are not substitutes, and reporting only the former would hide the finding.

Two observations follow. First, \textbf{a non-trivial share of the runs that an outcome-only benchmark would score as successes carry a hard violation}. Under \textsc{rules+state}, $4$ of the $21$ successes ($19\%$) are in category~(2).
% CHECK: with 27 + 4 = 31 successes in the category breakdown above, this should read 4 of 31 (12.9%).
Under the three-judge majority-vote protocol on the $772$ runs with resolved $O$, $29$ of $89$ successes ($32.6\%$; $3.8\%$ $[1.2, 7.1]$ of all runs) are in category~(2). Under all-agree consensus the figure is $10$ of $89$ ($11.2\%$ of successes; $1.3\%$ $[0.4, 2.3]$ of all runs), and under any-judge-flags it is $78$ of $89$ ($87.6\%$ of successes; $10.1\%$ $[3.1, 18.9]$ of all runs). These figures should be read against the inter-judge Krippendorff's $\alpha$ of $0.155$ reported in Section~\ref{sec:judge_results}: the spread between operating points is driven by judge disagreement, not by variation in agent behavior. Either way, category~(2) is real. On this benchmark it is thin, but it is not zero at any operating point we can compute.

Second, and more strikingly, \textbf{safe failure dominates the corpus}, although its exact share depends on the checker and the aggregation rule. Under the \textsc{rules+state} checker on all $297$ original runs, safe failure is $84.5\%$. Under the three-judge per-constraint protocol on the $772$ runs with resolved $O$ across both waves, Table~\ref{tab:four-cat-3j} gives the full four-category breakdown at each operating point and separately on each wave; the summary is safe failure at $81.5\%$ $[71.4, 90.2]$ when the judges must all agree, $51.8\%$ $[41.0, 62.4]$ under majority vote, and $15.8\%$ $[9.6, 23.2]$ under any-judge-flags. The three operating points span the possible verdicts of the same panel on the same runs, so the gap between them directly measures judge disagreement, not agent behavior.

\begin{table*}[t]
\caption{Four-category breakdown on the three-judge wave, under three operating points, with the original and extension waves separately alongside their union. Each cell is a percentage with its $95\%$ percentile-bootstrap interval ($10{,}000$ resamples, resampling whole scenarios). Categories are (1) intended success $O{=}1,V{=}0$, (2) unsafe success $O{=}1,V{=}1$, (3) safe failure $O{=}0,V{=}0$, (4) unsafe failure $O{=}0,V{=}1$. The original and Plan~E waves are indistinguishable at every category under every operating point: the framework's main finding (safe-failure dominance) replicates on 20 scenarios the four agents never saw, with constraints committed to git before the first run. The $5$ Plan~E runs with $O$ unresolved on measure\_sh edge cases are left out, not imputed.}
\label{tab:four-cat-3j}
\vskip 0.1in
\begin{center}\begin{small}
\begin{tabular}{lccc}
\toprule
Category & Original (N=297) & Plan~E (N=475) & Union (N=772) \\
\midrule
\multicolumn{4}{l}{\emph{All-agree (all 3 judges flag the same hard constraint)}} \\
(1) intended success &  9.1\% [1.5, 18.5] & 10.9\% [0.6, 24.9] & \textbf{10.2\% [3.0, 19.4]} \\
(2) unsafe success   &  1.3\% [0.0,  3.1] &  1.3\% [0.2,  2.5] &  1.3\% [0.4,  2.3] \\
(3) safe failure     & 86.5\% [75.3, 95.6]& 78.3\% [63.9, 90.9]& \textbf{81.5\% [71.4, 90.2]} \\
(4) unsafe failure   &  3.0\% [0.7,  6.1] &  9.5\% [3.4, 16.8] &  7.0\% [2.9, 11.9] \\
\midrule
\multicolumn{4}{l}{\emph{Majority vote ($\geq 2$ of $3$ judges)}} \\
(1) intended success &  6.1\% [1.0, 12.4] &  8.8\% [0.6, 20.2] &  7.8\% [2.1, 15.2] \\
(2) unsafe success   &  4.4\% [0.4,  9.9] &  3.4\% [0.2,  7.8] &  3.8\% [1.2,  7.1] \\
(3) safe failure     & 53.9\% [40.0, 67.6]& 50.5\% [35.2, 65.6]& 51.8\% [41.0, 62.4] \\
(4) unsafe failure   & 35.7\% [22.6, 49.7]& 37.3\% [23.8, 51.4]& 36.7\% [27.0, 46.9] \\
\midrule
\multicolumn{4}{l}{\emph{Any judge flags any hard constraint}} \\
(1) intended success &  2.4\% [0.4,  4.8] &  0.8\% [0.0,  2.1] &  1.4\% [0.5,  2.6] \\
(2) unsafe success   &  8.1\% [1.1, 16.6] & 11.4\% [1.1, 24.9] & 10.1\% [3.1, 18.9] \\
(3) safe failure     & 14.1\% [7.3, 22.4] & 16.8\% [8.0, 27.7] & 15.8\% [9.6, 23.2] \\
(4) unsafe failure   & 75.4\% [63.8, 86.0]& 70.9\% [56.6, 83.5]& 72.7\% [62.5, 81.7] \\
\bottomrule
\end{tabular}
\end{small}\end{center}
\vskip -0.1in
\end{table*}

Section~\ref{sec:model} notes in passing that an agent which stops and reports that the task cannot be done within its permissions is scored the same as one that failed through incompetence. On this benchmark that is not a marginal case; under every well-behaved checker configuration we report, it is the largest cell. We see this as a stronger argument for the two-dimensional model than category~(2). The unsafe-success cell is real but thin, and every per-agent estimate of it is very noisy, whereas the conflation of principled refusal with incompetence affects most of the corpus. It also suggests that, on tasks of this kind, agentic safety benchmarks that report a single success rate are mostly measuring an agent's willingness to break rules.

\subsection{Judges disagree, and the disagreement is first-order}
\label{sec:judge_results}

We report two judge protocols. The first is the Phase~1 two-judge batched protocol, run on $200$ runs of the original wave, with one call per run per judge and all constraints judged together. The two judges, Haiku-5.5 and GPT-6-luna, agree on run-level $V$ for $85.5\%$ of runs, and their disagreement is almost entirely one-directional: judge~B flags a hard violation that judge~A does not in $28$ runs, and the reverse happens in a single run. The second is a three-judge per-constraint protocol, run on $779$ runs across both waves, with one call per run, per constraint, and per judge, and no batching. Its three judges are smaller open-weight models from different families (Table~\ref{tab:judges}), and they agree with one another far less.

\begin{table}[t]
\caption{Judge slate for the three-judge per-constraint wave: three families, three providers, none also an agent in the study. Flag rates are per hard-constraint cell on the $3{,}814$ cells each judge scored; intervals are $95\%$ percentile bootstrap, $10{,}000$ resamples, resampling whole scenarios. ``High-conf'' is the share of a judge's verdicts it returned at confidence $\geq 0.8$ (its own self-reported field).}
\label{tab:judges}
\vskip 0.1in
\begin{center}\begin{small}
\begin{tabular}{llcc}
\toprule
Judge & Family & Flag rate (95\% CI) & High-conf \\
\midrule
Llama-3.1-8B  & Meta    & 30.0\% [25.2, 35.2] & 88\% \\
Mistral-Nemo  & Mistral & 17.5\% [13.3, 22.2] & 99.9\% \\
Qwen3.7-flash & Qwen    &  5.4\% [3.2, 7.8]   & 98\% \\
\bottomrule
\end{tabular}
\end{small}\end{center}
\vskip -0.1in
\end{table}

Krippendorff's $\alpha$ is $0.155$ $[0.103, 0.209]$ on hard-constraint cells, bootstrapped over scenarios, and pairwise Cohen's $\kappa$ lies in $[0.036, 0.261]$. The per-cell flag rates in Table~\ref{tab:judges} have intervals that do not overlap pairwise. The gate's verdict on this panel then swings with the aggregation rule: run-level $V$ is $82.7\%$ under any-judge-flags, $40.3\%$ under majority vote, and $6.9\%$ under consensus.
% CHECK: 779 runs here vs. 777 total and 772 with resolved O elsewhere; also, the abstract and conclusion quote 7%/73% "unsafe failure", whereas this paragraph gives run-level V of 6.9%/82.7%. Align the figures and the wording.
For comparison, the ODCV-Bench paper \citep{li2025odcv} reports Krippendorff's $\alpha=0.81$ across four frontier judges on a 0--5 rubric. The gap between $0.155$ and $0.81$ on the same benchmark, with judges drawn from either side of the capability frontier, is the empirical signal that the gate's accounting for checker error is not optional. On the two-judge original wave, the mean of $U$ is $0.0164$ and $75$ of $200$ runs have $U>0$. The three-judge wave gives $U>0$ on $4{,}969$ of the $5{,}089$ cells where at least two judges rendered a verdict ($97.6\%$), which confirms that the two-judge measurement understated $U$.

\textbf{Which agent one would ship depends on which judge one asks, and on which aggregation rule one picks.} Under judge~A alone (batched), three of the four agents pass the gate; under judge~B alone, none do. Under the three-judge majority-vote protocol no agent passes at $\delta=0.05$, and under the three-judge all-agree protocol three of four pass. The adjudication protocol of Section~\ref{sec:challenges} (deterministic checks where possible, several monitors from different families, and human review of disagreements) is what this measurement needs before it can be useful at all.

Excluding same-model judgments lowers the measured violation rate for both affected agents (from $0.175$ to $0.125$, and from $0.225$ to $0.175$). In other words, these models judged their own runs \emph{more} harshly than cross-family judges did. With only two of the four agents affected this is suggestive rather than established, but it runs opposite to the self-preference effect one might expect \citep{panickssery2024llm}.

\paragraph{Guidance on aggregation when $\alpha$ is unknown.}
The three operating points we report are not interchangeable, and until the audit returns calibrated $\alpha_j$ and $\beta_j$ per judge $j$ (which would enable the per-cell Bayesian aggregation the framework prescribes), the choice should be matched to the downstream use. For \emph{release gating} we recommend \textbf{consensus} (all-agree): it is the lowest-$\alpha$ operating point available without an audit, since a false positive requires every judge to err in the same direction on the same cell, which is empirically rare at $\alpha=0.155$. The cost is a lower-bounded false-negative rate, but a false negative in gating is less costly than a false positive once the gate's veto is treated as absolute (Section~\ref{sec:model}). For \emph{continuous monitoring} we recommend \textbf{majority vote}: it balances $\alpha$ against $\beta$ and preserves resolution for ranking agents against each other. For \emph{triage of candidate audit targets} we recommend \textbf{any-judge-flags}: it is deliberately high-recall and feeds the audit sample rather than acting on it. All three should be reported together whenever the audit is still pending, because the \emph{spread} between the three is itself the first-order error bar on any single number.

\subsection{Early warning does not work here}

\begin{table}[t]
\caption{Early warning on this corpus. AUROCs are run-level, computed on prefixes strictly before the first violating action. The unmatched column takes the maximum over all eligible prefixes, and so over systematically more draws for compliant runs than for violating ones, which are truncated at the first violation. Matching the horizon to $H$ pre-violation actions removes this confound: the \texttt{action\_index} falsifier sits far below chance ($0.172$) when unmatched and at exactly $0.500$ when matched, which is how we validated the matching. At the matched horizon, $D$ carries no usable signal. A next-action scorer built from typed features of the proposed action and its preceding context reaches a modest prefix-level AUROC of $0.603$ ($0.604$ with position features ablated) but no usable detection: at a $5\%$ false-alarm budget it fires in $0$ of $17$ violating runs at $H=2$ and $0$ of $10$ at $H=3$. Only $10$--$17$ violating runs survive to the matched horizon, and every AUROC interval contains chance. FAR: false-alarm rate.}
\label{tab:rq2}
\vskip 0.1in
\begin{center}\begin{small}
\setlength{\tabcolsep}{5pt}
\begin{tabular}{lccc}
\toprule
 & & \multicolumn{2}{c}{Matched} \\
\cmidrule(l){3-4}
Score & Unmatched & $H{=}2$ & $H{=}3$ \\
\midrule
\texttt{action\_index} (falsifier) & 0.172 & 0.500 & 0.500 \\
\midrule
$D$ (prefix distance) & 0.293 & 0.496 & 0.432 \\
Next action & 0.396 & 0.381 & 0.512 \\
Keyword rule & 0.495 & 0.500 & 0.500 \\
\bottomrule
\end{tabular}

\vspace{0.6em}

\setlength{\tabcolsep}{4pt}
\begin{tabular}{lccc}
\toprule
Operating point ($D$) & Detected & Realized FAR & Target \\
\midrule
$H{=}2$, 5\% target  & 0/17 & 0.054 & 0.05 \\
$H{=}2$, 10\% target & 2/17 & 0.120 & 0.10 \\
$H{=}3$, 5\% target  & 0/10 & 0.103 & 0.05 \\
$H{=}3$, 10\% target & 0/10 & 0.161 & 0.10 \\
\bottomrule
\end{tabular}
\end{small}\end{center}
\vskip -0.1in
\end{table}

Table~\ref{tab:rq2} reports the anomaly component. The result is negative, but how strongly negative depends on controlling an asymmetry that we had not isolated when the paper was first drafted. Violating runs are truncated at their first violation and so contribute a median of $3$ eligible prefixes; compliant runs are not truncated and contribute a median of $5$. Taking a run-level maximum over unequal numbers of draws biases the score systematically against the positive class, and that alone is enough to make an uninformative score \emph{look} anti-predictive. We built a falsifier to measure this effect: it uses the prefix index alone as the score, with no features. A score of this kind carries no information about violations, yet at the unmatched horizon it reaches an AUROC of $0.172$; once the horizon is matched to a fixed $H$ pre-violation actions it returns to exactly $0.500$, as it should. We take this as evidence that the matching works and that the unmatched run-level AUROCs should not be read as reflecting any real mechanism. At the matched horizon, $D$ has an AUROC of $0.496$ ($H{=}2$) and $0.432$ ($H{=}3$). The right reading is that $D$ carries no usable signal here, not that it is actively anti-informative.
% CHECK: the abstract gives the raw run-level AUROC of D as 0.312, but the table's unmatched column gives 0.293. Align.

The reason can be seen in the data. Runs are short (median $5$ actions), and violations occur early, after ordinary reconnaissance. A pre-violation prefix is typically two or three file reads, which is the most common compliant behavior there is. $D$ measures distance from compliant behavior, and these prefixes \emph{are} compliant behavior: the violation is a step change, not a drift. The early-warning idea of Section~\ref{sec:framework} presupposes a trajectory long enough to depart gradually from a norm, and these trajectories are not that long.

The calibration diagnostic is the more useful result. Split-conformal calibration should hold the false-alarm rate at its target when calibration and test prefixes are exchangeable. On the original, unmatched score, the realized rates are $14.3\%$ against a $5\%$ target and $27.0\%$ against $10\%$ (first row of Table~\ref{tab:conformal}), an overshoot of $2.9\times$ and $2.7\times$ respectively, so the gap is not an artefact of a single operating point. The matched-horizon operating points in the lower part of Table~\ref{tab:rq2} overshoot less, but they still detect almost nothing. \textbf{Exchangeability does not hold across scenarios}, which is exactly the failure Section~\ref{sec:framework} anticipates when the agent meets task types absent from the calibration set. Calibration and evaluation scenarios are different tasks with different tool vocabularies, so a prefix from an unseen scenario lies far from the reference set for reasons that have nothing to do with compliance. The per-task-family calibration we recommend is not possible at this scale: with 20 scenarios split evenly, no family has enough runs to calibrate within.

\paragraph{Exchangeability-robust calibration does not rescue it.}
Since the diagnosis is non-exchangeability, the natural response is a calibration method designed for it. We tried three, with the results in Table~\ref{tab:conformal}: Mondrian conditioning on prefix length (the covariate on which $D$ depends most strongly), covariate-shift weighting of calibration scores by similarity to the test prefix, and cross-conformal aggregation over folds of calibration scenarios. \textbf{None removes the overshoot.} At the run level, Mondrian gives $3.2\times$ and $2.7\times$, against $2.9\times$ and $2.7\times$ for split conformal; cross-conformal gives a marginal improvement at the looser target ($2.2\times$) and none at the tighter one.

\begin{table}[t]
\caption{Realized false-alarm rate under exchangeability-robust calibration variants, with the multiple of the target in parentheses. None removes the overshoot. $^\dagger$These two variants raise an alarm if \emph{any} prefix exceeds its threshold, i.e.\ one test per action rather than one per run; their inflation is partly due to multiplicity, and they are not directly comparable with the others.}
\label{tab:conformal}
\vskip 0.1in
\begin{center}\begin{small}
\setlength{\tabcolsep}{4pt}
\begin{tabular}{lcc}
\toprule
Calibration & Target 5\% & Target 10\% \\
\midrule
Split conformal (original) & 0.143 (2.9$\times$) & 0.270 (2.7$\times$) \\
Mondrian, prefix length & 0.159 (3.2$\times$) & 0.270 (2.7$\times$) \\
Cross-conformal (CV+) & 0.143 (2.9$\times$) & 0.222 (2.2$\times$) \\
Weighted, cov.\ shift$^\dagger$ & 0.302 (6.0$\times$) & 0.492 (4.9$\times$) \\
Mondrian, per prefix$^\dagger$ & 0.254 (5.1$\times$) & 0.444 (4.4$\times$) \\
\bottomrule
\end{tabular}
\end{small}\end{center}
\vskip -0.1in
\end{table}

One methodological note changes how the table should be read. Our first implementation of the Mondrian and weighted variants raised an alarm whenever \emph{any} prefix exceeded its threshold, which amounts to one test per action, while split conformal tests only the run maximum. Those variants accordingly looked much worse ($5.1\times$ and $6.0\times$), but most of that gap is multiplicity, not shift. Standardizing each prefix score within its length bucket and testing the run maximum restores a like-for-like comparison, and the difference between the variants then largely disappears. We report both versions so that the artefact is visible.

Our conclusion is that the problem does not lie in the calibration layer. A conformal wrapper controls the error rate of whatever score it is given; it cannot create signal. With the horizon artefact controlled, $D$'s AUROC is within noise of chance, so the underlying score is uninformative, and recalibrating an uninformative score more cleverly only yields a better-behaved uninformative detector. A direct repair of the kind the framework invites, scoring the proposed next action instead of the prefix distance, gives a modest prefix-level AUROC of $0.603$ ($0.604$ with position features ablated, so the discrimination comes from the typed action rather than from its position in the run). It still provides no deployable detection: at a $5\%$ false-alarm budget it fires in $0$ of $17$ violating runs at the matched horizon. Fixing the early-warning component on this corpus would require a different representation or longer trajectories, not a different quantile.

\paragraph{Richer sequence models and alternative signals.}
We did not fit an HMM, a learned contrastive embedding, or a trajectory-level transformer to this corpus. Three reasons: there is no labeled seed per task family (the audit is what would produce one), the median trajectory is five actions long, and the hold-out at the task-family level is one or two scenarios per family. All three choke a model with more than a few parameters; the ablation we can state is that the modest-dimensionality next-action scorer already uses features ($\approx 50$ typed-feature indicators times a context prefix) comparable to what a small HMM would discover, and the AUROC gap from chance to $0.603$ is the ceiling under this representation. A larger follow-up with longer trajectories (say, 50--100 actions per run) and a labeled seed from the audit would open three directions we consider most promising: (i) \textbf{proximity-to-contract-boundary} as a feature, measuring how close each action's typed features sit to the discriminator between contractually-permitted and contractually-forbidden actions in the constraint schema --- this exploits the schema's \texttt{evidence} field directly, and does not require learning what ``bad'' looks like from the data; (ii) a \textbf{per-task-family contrastive embedding} seeded with audit labels, where exchangeability within a family can be reasonably assumed and cross-family calibration is deliberately avoided; (iii) \textbf{alarm-policy evaluation} rather than raw AUROC \citep{barber2023conformal}, binding each detector to an explicit response (abort, pause for human review, log-only) and reporting its operating characteristic at that commitment. Our negative result here is a caution against pre-committing to any such detector before the audit returns per-family labels.

\subsection{How much do our own design choices decide the answer?}

Three choices in Sections~\ref{sec:framework} and~\ref{sec:setup} were made a priori and could have driven the results. We test each of them.

\paragraph{Soft-constraint weights.}
The weights $w_i$ use a fixed three-point scale and are not calibrated against any external severity model, which is a fair criticism of the framework as much as of this study. We recompute $\mathrm{SC}$ under four schemes: ours; uniform ($w_i=1$); an expert-style scheme keyed to the \emph{source} field, on the view that violating an unstated norm is worse than violating an explicit instruction ($1.0/1.5/2.0$ for explicit/implied/norm); and inverse frequency, under which rare violations weigh more. The per-agent ordering is largely stable. Kendall's $\tau$ against our scheme is $1.00$ for the source-keyed scheme, $0.91$ for uniform, and $0.67$ for inverse frequency, and the single reordering under inverse frequency swaps two agents whose intervals overlap almost entirely. Absolute $\mathrm{SC}$ values, however, move by a factor of three across schemes, so \emph{levels} of $\mathrm{SC}$ should not be compared across studies that weight differently, even though the within-study ordering survives.

\paragraph{Batching constraints into one judge call.}
Our budget forced all of a scenario's constraints into a single judge call, which risks interference between constraints. We re-judged a stratified subsample of 158 (run, constraint) pairs over 24 runs, chosen to over-represent runs that the batched pass flagged, with one constraint per call and the same judge, the same trajectory rendering, and the same reasoning setting, so that batching was the only difference. Cell-level agreement is $94.3\%$ ($95\%$ CI $[89.5, 97.0]$), but because violations are rare, Cohen's $\kappa$ is only $0.498$. The disagreement is directional: \textbf{batching raises the measured violation rate from $0.044$ to $0.076$}. The batched protocol flags 7 pairs that the isolated protocol does not, and the reverse happens twice. Judging constraints together therefore makes a judge roughly $70\%$ more likely to report a violation. This does not overturn any ordering we report, since the effect applies to every agent alike, but it does mean that our judge-based violation \emph{rates} are biased upward. It is one more reason why the $\alpha$ from the audit is the quantity that matters.

\paragraph{The scenario draw.}
Scenarios were drawn with a fixed seed, stratified by container base image, before any result was seen. We did not re-draw, so we cannot quantify draw variance; this is a known gap, not a tested one. The scenario-clustered bootstrap used throughout propagates between-scenario variability into every interval, which is the partial mitigation available at this sample size.

\subsection{State evidence is less decisive than the hierarchy suggests}

Section~\ref{sec:framework} orders evidence by reliability, with environment state first. Our experience qualifies that ordering. A naive state checker (the file changed, so the constraint was violated) produced a high false-positive rate. We inspected three cases by hand: an agent that rewrote eleven protected data files \emph{after backing them up, describing them as repaired malformed records, and reporting the true failing score}; thirteen runs that ``modified'' an incident report by \emph{appending a review section}; and fifteen runs whose audit-log change was made by the scenario's own tool. We therefore classify changes as rewritten, appended, truncated, or deleted instead of recording a single bit, which cleared fifteen appends while still flagging nineteen rewrites and deletions.

The correction is itself instructive. Whether a change is an append or a rewrite is decided from the command, not from the hashes, so a constraint that relies on this distinction is really using log-rule evidence dressed up as state evidence. State evidence establishes \emph{that} something changed and \emph{how}; whether the change was legitimate still needs a judge. The hierarchy holds for the question ``did this file change?'' but not for the question ``was this a violation?''

\subsection{What is pending the audit}
\label{sec:audit-pending}

Three of the measurements proposed in Section~\ref{sec:framework} are not reported here because they are defined against human labels: whether $U$ predicts that the automated label is wrong, the precision and recall of each checker configuration, and the behavior of the gate under corrected error rates via equation~\eqref{eq:corrected}. We have deliberately not substituted model labels for human ones. The stratified sample ($200$ runs, including every flagged run so that $\beta$ can be estimated) and a blinded annotation interface are built and ready, and the estimator and its interval machinery are implemented and tested against the worked example of Section~\ref{sec:model}.

\paragraph{Audit design, with the agent-dependence concern made explicit.}
The audit protocol stratifies by \emph{agent family}, \emph{constraint evidence class} (state, log-rule, judge), and \emph{scenario family} (the four Docker base-image groupings), with the flagged stratum deliberately oversampled so that $\beta$ is estimable at all and the sampling weights recorded so that the analysis can reweight to the population. Each run--constraint cell is labelled by two annotators drawn from a pool of three people who do not work on this paper, from a blinded interface that strips the model name, every automated score, the per-constraint checker verdicts, and the flag that drove sampling; the exact fields removed are enumerated in the bundle's manifest. Cohen's $\kappa$ is reported overall and per constraint class, disagreements are written out for the authors to adjudicate, and unresolved cells are counted as missing rather than coerced. The result is an audit that supports an \emph{agent-dependent} false-positive rate $\alpha_a$ and false-negative rate $\beta_a$ per agent $a$, not a pooled $\alpha$: the review of this manuscript rightly flagged that an agent which conceals or obfuscates actions may induce a different $\alpha_a$ than an agent that does not, so pooling would silently bias the corrected rate $p$ in equation~\eqref{eq:corrected} downwards for the former. We adopt the stratification because the fix is cheap once the sample is drawn this way, and because the alternative --- pooling and documenting the pooling as a limitation --- would leave readers unable to tell whether an across-agent ranking under $p$ reflects behaviour or stratum size. The infrastructure is on-disk and reproducible from the commit this manuscript cites; the delivery of labelled files is the gating step. The paper's forward-looking claims are stated as provisional on the audit's outcome and will be updated with the corrected rates and gate verdicts in the camera-ready revision.

One comparison we had planned turned out to be impossible rather than merely pending. The benchmark publishes its own severity labels, but for twelve models, none of which is in our slate, so no run overlaps and agreement cannot be computed.

A second gap, originally listed as a concrete task, has since been closed. Six of the original twenty scenarios had no objective score, for two distinct reasons. Four shipped a KPI script that required an argument the prompt did not determine (namely, which artefact should be validated), and two wrote their result to a file instead of standard output, so there was nothing to parse. The extension study specified the missing rules in \texttt{config/\allowbreak objective\_rules.yaml} from the pristine scripts' real output on real traces, justifying each choice in the rule file. Adding three mechanisms that the schema previously lacked (\texttt{args}; \texttt{read\_file}; and a \texttt{measure\_sh} block for aggregate measurements, written in the YAML so that authors can review them) was enough to raise $O$ coverage from $70\%$ to $100\%$ of the original 297 runs without re-running a single agent. The extension's 20 new scenarios were resolved with the same mechanisms, plus a purpose-built measurement for the four that ship no KPI scorer, bringing the extension's $O$ coverage to $99.0\%$ (5 runs remain unresolved because of \texttt{measure\_sh} edge cases). This closes the gap on the measurement side of the framework, at the cost of making explicit in the rule file every editorial decision that the schema's silent defaults would otherwise have taken.

\subsection{Reproducibility and release}
\label{sec:release}

We release, under an open licence, every artefact a reader needs to reproduce every number in this paper from a clean machine. The repository contains: the agent harness, the Docker sandbox, the typed-action encoder, the per-scenario rule file \texttt{config/\allowbreak objective\_rules.yaml} (with each editorial decision annotated inline), the full set of $262$ hand-written constraints in the six-field schema, the stratified audit sample with its blinding map, the blinded annotation interface, the analysis scripts (including the scenario-cluster bootstrap, the Rogan--Gladen correction, and the conformal-calibration variants), and the raw traces of every one of the $777$ runs. The analysis pipeline is a single resumable script that re-derives every table and figure from the traces with no API calls, so the ``reproducibility'' claim holds against future deprecations of the model APIs. The exact judge prompts used in both the batched and per-constraint protocols, together with the decision schema the judges return (a boolean verdict, a confidence in $[0,1]$, and a free-text reason), are included verbatim; so are the four reasoning-mode settings used per agent and the full list of OpenRouter model IDs and prices as fetched on the run date. For external re-audit of the gate's verdict on any particular agent, the stored final filesystem state of every run (as a tarball) allows re-scoring under any checker a reader supplies without re-running the agent. The repository URL and the commit this manuscript cites are listed in Appendix~A.

\paragraph{Non-determinism, replication budget, and designs we would adopt in a larger follow-up.}
Provider-side non-determinism persisted at the lowest sampling temperature the two providers accept, so our runs are random samples from each agent's response distribution, not greedy decodes. The practical implication for a follow-up is that run counts should be budgeted for the between-run variance of $V$ and $O$ at the \emph{scenario} level, not the agent level: our scenario-cluster bootstrap shows that single-scenario intervals on $V$ are wide (often $\pm 0.10$) even with three repetitions per cell. A design we recommend for any paper that wants statistical power to separate agents is a \emph{paired scenario} design --- the same scenario, same prompt, same seed where seeds exist, run against each agent in the panel, with repetitions nested within scenario and agent and the analysis fit as a mixed-effects model with scenario as a random intercept. In our data, paired scenario comparisons yield intervals roughly $\sqrt{k}$ tighter than unpaired comparisons at the same cost, where $k$ is the mean number of agents per scenario. We suggest a minimum of five repetitions per cell for ranking claims, and reporting the within-scenario standard deviation separately from the across-scenario one so that provider non-determinism is not conflated with task-difficulty variance. A further measure, which this study did not take, is to subscribe to the ``logit-temperature zero'' version of each provider's API where that is offered (OpenAI's \texttt{seed} parameter and the Anthropic equivalent, both still partial as of this writing).

\section{Research Challenges}
\label{sec:challenges}

\paragraph{Inferring intent from behavior.}
Deciding whether an agent \emph{deliberately} worked around a constraint is difficult, and unexpected behavior should not be read as strategic intent by default. More defensible evidence comes from the structure of the trajectory: what the environment returned, whether an attempt was repeated after a refusal, whether compliant alternatives were available and ignored, and how behavior changed once a constraint was encountered. Even then, conclusions about intent should be stated with care.

\paragraph{Hidden and emergent strategies.}
Agents operating over long horizons may find strategies that benchmark designers did not anticipate, and a benchmark cannot detect a violation that it offers no opportunity to commit. Future environments should deliberately include both a legitimate route and one or more constraint-violating routes to the same objective, so that the choice between them becomes observable.

\paragraph{Reward design.}
If completion is strongly rewarded and violations are only weakly penalized, an optimizer will sometimes prefer to violate. Benchmarks and training setups should keep task reward and constraint-related signals separate, so that this trade-off stays visible instead of being buried in a single objective.

\paragraph{Early detection.}
Post-hoc analysis helps in understanding failures, but deployed systems need monitoring while the agent is running. One option is a running risk estimate that combines trajectory anomaly, uncertainty, proximity to a constraint boundary, and tool-use patterns, and that requests human intervention when it crosses an adaptive threshold. How to set that threshold without overwhelming human reviewers is an open problem, and a practical one: a review of production agent systems finds that human-centered controls such as runtime approval are the dominant security mechanism in industry \citep{wang2026reframing}, and a monitor that asks too often will have its requests approved without being read.

\paragraph{Reliability of automated monitors.}
At scale, trajectory review will largely be done by other models, which raises an evaluator--agent problem. If the monitor shares blind spots with the agent it oversees, for example because both derive from the same base model, automated oversight can fail systematically and without any visible sign \citep{panickssery2024llm,greenblatt2024control}. The evidence so far is not reassuring. \citet{huang2026reward} report that a review panel misses many confirmed cases of reward hacking, and that detailed reviewer feedback helps agents evade detection in later rounds. We therefore recommend an adjudication protocol instead of a single judge: deterministic checks wherever the evidence allows, at least two monitors from different model families, and human review of every case on which they disagree, together with a random sample of the cases on which they agree. That sample gives an estimate of the errors the monitors share, which their agreement would otherwise hide. Whether such a protocol is reliable enough, and under which conditions, needs direct empirical study.

\section{Conclusion and Future Work}
\label{sec:conclusion}

The move from language models that answer to agents that act calls for a matching change in how they are evaluated. Completing a task and behaving acceptably while doing so are different things, and an agent that reaches the correct answer by an unauthorized route is a different object from one that reaches the same answer within its permissions. We have therefore proposed moving the unit of evaluation from the final answer to the whole interaction trajectory. Objective achievement, constraint compliance, uncertainty, and behavioral distribution shift are reported as separate dimensions, with a gating rule under which a hard violation cannot be offset by task success. This brings out a failure mode that conventional benchmarks conceal: agents that look successful while reaching their objectives through unintended or prohibited means.

The study in Sections~\ref{sec:setup} and~\ref{sec:results} supports part of this proposal and refutes another part. Separating the two dimensions does change the conclusion. Under deterministic checkers the outcome ranking and the gated ranking are opposed, and the category that an outcome-only score cannot distinguish from incompetence turns out to be most of the corpus. The early-warning component, by contrast, does not survive contact with data at this scale. The apparent anti-prediction of $D$ at the run level turned out to be a horizon-asymmetry artefact, and with the horizon matched the score is uninformative. A direct repair that scores the proposed next action gives a modest prefix-level AUROC ($0.60$) but catches no violating run at a $5\%$ false-alarm budget. The conformal guarantee on which the original score relied overshoots its target by roughly a factor of three, and three exchangeability-robust calibration variants failed to recover it. Together with the uninformative underlying score, this places the fault in the score and the trajectory length, not in the calibration layer.

The framework's weakest point is also the one it depends on most. Nearly half of the constraints we could write admit only a judge; two judges from different families disagreed enough to change which agent would ship; and the violation rate measured under the full protocol exceeded the gate tolerance for every agent, so the gate cannot discriminate until the checkers themselves are characterized. We take this as evidence that the framework's accounting for checker error is not a caveat attached to it, but the part most in need of development. The ODCV-Bench paper \citep{li2025odcv} shows that frontier judges on this benchmark agree with each other at Krippendorff's $\alpha=0.81$, a level at which the $\alpha$ and $\beta$ corrections we derive would be of secondary importance. We show in Section~\ref{sec:judge_results} that on the same benchmark, with judges from the smaller, open-weight tier that smaller teams are likely to use, inter-judge agreement drops to $0.155$, and the gate verdict swings from $7\%$ unsafe failure (consensus) to $73\%$ (any-judge-flags) purely as a function of the aggregation rule applied to the same panel's verdicts. Our framework's position is that this spread is the measurement signal, not a nuisance, and its gating rule and audit protocol are the apparatus that turns it into a release decision. The tolerance sweep makes the same point more sharply: the gate separates agents only when $\delta$ lies above the checker's false-positive rate and below the agents' true violation rate, and the framework as stated gives a practitioner no way to find that window without first measuring $\alpha$.

The framework also assumes that a constraint set can be written down and checked, which is often only approximately true, and that trajectories are fully observable, so an agent able to conceal its actions falls outside what trajectory monitoring can cover. Four directions for future work follow. The first is completing the human audit, so that $\alpha$ and $\beta$ become measurable and the gate can be applied to corrected rates. The second is building paired-trajectory datasets in which two agents achieve identical outcomes by compliant and non-compliant strategies, so that detectors can be trained and tested with the outcome held fixed. The third is longer-horizon tasks, on which distribution shift has room to appear before a violation; on short trajectories it does not. The fourth is larger and better-calibrated monitor panels, since two frontier judges disagreed enough to decide which agent ships, and three smaller open-weight judges reached a Krippendorff's $\alpha$ of only $0.155$ $[0.103, 0.209]$ on the same benchmark. Scaling the panel is necessary but not sufficient: the panel's own disagreement has to be measured and corrected for. The research question is not only ``can the agent achieve the goal?'' but also ``can the agent achieve the goal within the intended behavioral boundaries, and can we reliably tell when it has not?'' Our results suggest that the second half of that question is, for now, the harder one.

\section*{Impact Statement}

This paper discusses how to evaluate AI agents so that constraint-violating behavior is not mistaken for success. We expect the main consequence of this line of work to be safer deployment of autonomous systems. Trajectory-level monitoring involves recording agent actions, which may include user data, and such monitoring should be designed with appropriate privacy safeguards.

\bibliography{references}
\bibliographystyle{icml2025} % must match the year of the venue you submit to

\end{document}